% Part: counterfactuals
% Chapter: introduction
% Section: counterfactuals

\documentclass[../../../include/open-logic-section]{subfiles}

\begin{document}

\olfileid{cnt}{int}{cnt}

\olsection{反事实条件句}

英语中一种非常常见且重要的“如果……那么……”结构，是使用 \emph{to be} 的过去虚拟语气构建的：“如果情况是……那么情况就会是……”。由于这类条件句的前件通常为假，即与事实相反，因此它们被称为\emph{反事实条件句}（又因它们使用了 \emph{to be} 的虚拟语气，也称为\emph{虚拟条件句}）。它们不同于\emph{直陈}条件句，后者的形式为“如果情况是……那么情况就是……”。反事实条件句与直陈条件句的真值条件不同。考虑 Adams（亚当斯）的著名例子：

\begin{quote}
  如果奥斯瓦尔德没有刺杀肯尼迪，那么是别人刺杀的。

  如果奥斯瓦尔德当时没有刺杀肯尼迪，那么也会有别人刺杀。
\end{quote}

第一句是直陈条件句，第二句是反事实条件句。第一句显然为真：我们知道约翰·F·肯尼迪总统是被\emph{某人}刺杀的，而如果那个人不是（与沃伦委员会报告结论相反）李·哈维·奥斯瓦尔德，那么就是别人刺杀了肯尼迪。第二句说的则是另一回事。它声称，如果奥斯瓦尔德没有刺杀肯尼迪，即如果达拉斯的枪击事件得以避免或未能成功，那么历史将以另一种方式展开，使得另一场刺杀得以成功。要使这句话为真，就必须有强大的势力曾密谋确保肯尼迪死亡（正如许多肯尼迪刺杀阴谋论者所相信的那样）。

关于\emph{直陈}条件句能否被实质条件句正确刻画，目前仍存在激烈争论，特别是关于实质条件句的悖论能否以某种与实质条件句给出英语直陈条件句真值条件相容的方式得到“解释”。相比之下，反事实条件句不能用实质条件句正确符号化，这是毫无争议的。原因很清楚：尽管反事实条件句的前件通常为假，但并非所有前件为假的反事实条件句都为真——例如，如果你相信沃伦委员会报告，并且不存在刺杀肯尼迪的阴谋，那么 Adams 的那个反事实条件句就是一个反例。

反事实条件句在因果推理中扮演着重要角色：反事实条件句的一个主要用途就是表达因果关系。例如，划火柴会导致它点燃，你可以通过说“如果这根火柴被划了，它就会点燃”来表达这一点。实质条件句，以及一般的直陈条件句，不能用来表达这一点：如果火柴从未被划过，那么“火柴被划了 $\lif$ 火柴点燃了”就为真，无论假如它被划会发生什么。更糟的是，如果火柴从未被划过，“火柴被划了 $\lif$ 火柴变成了一束花”也为真，但火柴倘若被划，显然不会变成一束花。

反事实条件句的正确逻辑究竟是什么，至今仍有争论。Stalnaker（斯塔尔内克）和 Lewis（刘易斯）给出了一种有影响力的反事实条件句分析。根据他们的分析，一个反事实条件句“如果 $S$ 成立，那么 $T$ 就会成立”为真，当且仅当在与实际情况最接近且 $S$ 为真的反事实情境（“可能世界”）中 $T$ 为真。这被称为“本体论”分析，因为它涉及可能世界的本体论。其他分析则利用条件概率或信念修正理论。目前涌现了大量不同的反事实逻辑方案。甚至不存在一个单一的 Lewis–Stalnaker 反事实逻辑：尽管 Stalnaker 和 Lewis 沿着相似的思路、参照最接近的可能世界提出了各自的理论，但他们所做的假设导致了不同的有效推理模式。

\end{document}